\section{Fourier Transform}

\begin{equation*}
    \mathcal{F}\left\{x(t)\right\} \equiv X(\omega) \equiv \int_{-\infty}^{\infty} x(t) e^{-i\omega t} dt
\end{equation*}

The \textbf{Fourier transform} $\mathcal{F}$ maps a signal from the time domain to the frequency ($\omega$) domain. It can also be used on spatial data. A key advantage of the Fourier transform is data compression: frequency values often take up much less storage space than a full representation of the signal in the time domain.

\vsp
The convention is to use a lowercase function like $x(t)$ in the time domain and an uppercase function like $X(\omega)$ in the frequency domain.

\vsp
The integral usually yields a real and an imaginary component. The imaginary component contains phase information. However, we are usually only interested in the real component since that yields frequency information.

\subsection{Discrete Fourier Transform}

\begin{equation*}
    \mathcal{F}\left\{x(n)\right\} \equiv X(k) \equiv \sum_{n = 0}^{N-1} x(n) e^{-i\frac{2\pi kn}{N}}
\end{equation*}

where $N$ is the number of data points.

\vsp
The \textbf{discrete Fourier transform (DFT)} is suitable for computers. However, in practice, software uses the FFT.

\subsection{Fast Fourier Transform}

If $N = 2^q$ for some integer $q$, then the DFT can be computed much faster using the \textbf{fast Fourier transform (FFT)}. Whereas the regular Fourier transform and DFT take $O(N^2)$ time to compute, the FFT takes $O(N \log N)$ time to compute. Any number of data points can be brought to the next highest power of 2 by padding the ends with zeros (a technique called \textbf{zero-padding}).

\subsection{Fourier Transform Table}

There are many Fourier transform pairs, but these are the ones we will need in this class.

\begin{table}[h]
    \centering
    \begin{tabular}{cc}
        $x(t)$              & $X(\omega)$                                                    \\
        \hline
        \hline
        1                   & $2\pi \delta(\omega)$                                          \\
        $\delta(t)$         & 1                                                              \\
        $\rect(t/T)$        & $T ~ \sinc(\frac{T}{2\pi} \omega)$                             \\
        $\sinc(\omega t)$   & $\frac{1}{\omega_0} \rect(\frac{\omega}{2\pi \omega_0})$       \\
        $\cos (\omega_0 t)$ & $\pi[\delta(\omega - \omega_0) + \delta(\omega + \omega_0)]$   \\
        $\sin (\omega_0 t)$ & $-i\pi[\delta(\omega - \omega_0) - \delta(\omega + \omega_0)]$ \\
        \hline
    \end{tabular}
\end{table}

\subsection{Convolution Theorem}

\begin{equation*}
    x(t) \ast h(t) = \mathcal{F}^{-1}\left\{\mathcal{F}\left\{x(t)\right\} \mathcal{F}\left\{h(t)\right\}\right\}
\end{equation*}
Notice that multiplication in the time domain is equivalent to convolution in the frequency domain, and convolution in the time domain is equivalent to multiplication in the frequency domain.

This theorem is useful in because convolutions are much slower to compute than the Fourier transform.

\subsection{PSD}

[work in progress - this topic will be covered later in the class]

\begin{equation*}
    \text{PSD}(\omega) = \left| X(\omega) \right|^2 = (\text{real})^2
\end{equation*}

\subsection{Windowing Functions}

The Fourier transform inherently messes with the spectral content of a signal. To mitigate this, we can use windowing functions.

\begin{equation*}
    \mathcal{F}\left\{x(t)\right\} = \mathcal{F}\left\{x(t)w(t)\right\} = X(\omega) \ast W(\omega)
\end{equation*}

where $w(t)$ is the windowing function.

The perfect window needs an infinite number of data points, which is impossible to achieve in practice. Instead, we want to minimize the \textbf{width} of the \textbf{main lobe} in a dB-vs-frequency plot of the window function.

\subsubsection{Hanning Window}

\begin{equation*}
    W(n) = \frac{1}{2}\left[1 - \cos\left(\frac{2\pi n}{N-1}\right)\right]
\end{equation*}

PSD has multiple full-forms:

\begin{itemize}
    \item Power Spectrum
    \item Power Spectrum Distribution
    \item Power Spectral Distribution
    \item Power Spectral Density
\end{itemize}

Each are technically defined differently, but are often used interchangeably. We will stick to PSD to avoid any confusion.